\section{Framework}
\label{sec:framework}

The docstring--test--code closure framework consists of three components:
(i)~the \emph{docstring--test--code triangle} of artefacts associated with a
Python function, (ii)~three \emph{closure paths} that each reconstruct one
artefact from the other two via a sequence of two SLM stages, and
(iii)~a \emph{strict-AND validity policy} that pairs an SE-validated
automated metric with an external judge SLM's semantic-equivalence rating
to decide whether a reconstruction counts as closed. Figure~\ref{fig:framework}
gives the overall picture; the remainder of this section formalises each
component. The experimental instantiation of the framework (which SLMs sit
at which stages, which benchmark functions, and how statistical inference
is drawn) is deferred to Section~\ref{sec:experiment_setup}.

% Framework figure — three closure paths on the docstring–test–code triangle.
% Rendered natively in TikZ; no external image dependency.
\begin{figure*}[t]
  \centering
  \begin{tikzpicture}[
    node distance=1.4cm and 1.1cm,
    orig/.style   ={rectangle, draw, thick, minimum height=0.9cm, minimum width=1.0cm, font=\bfseries, fill=gray!8},
    recon/.style  ={rectangle, draw, thick, minimum height=0.9cm, minimum width=1.0cm, font=\bfseries, fill=blue!8},
    slm/.style    ={-{Latex[length=2.2mm]}, thick},
    metric/.style ={rectangle, draw=blue!60!black, thick, rounded corners=2pt,
                    minimum height=0.9cm, text width=3.2cm, align=center,
                    font=\footnotesize, fill=blue!5},
    judge/.style  ={rectangle, draw=orange!70!black, thick, rounded corners=2pt,
                    minimum height=0.9cm, text width=3.2cm, align=center,
                    font=\footnotesize, fill=orange!8},
    pathlabel/.style={font=\bfseries\small},
  ]

    %% ------------------- Path 1: C -> D -> T -------------------
    \node[orig]  (P1C)  at (0,3)     {$C$};
    \node[recon,right=of P1C] (P1D)  {$D'$};
    \node[recon,right=of P1D] (P1T)  {$T'$};
    \node[metric,right=of P1T] (P1M)
      {mutation kill rate\\$K(T', C)$};
    \node[judge, right=of P1M] (P1J)
      {judge SLM rating\\$J(D, D')$};

    \draw[slm] (P1C) -- node[above, font=\scriptsize]{$L_{\mathsf{spec}}$}  (P1D);
    \draw[slm] (P1D) -- node[above, font=\scriptsize]{$L_{\mathsf{test}}$}  (P1T);
    \draw[slm] (P1T) -- (P1M);
    \draw[slm] (P1M) -- (P1J);
    \node[pathlabel, left=0.15cm of P1C]{Path 1};

    %% ------------------- Path 2: D -> T -> C -------------------
    \node[orig,  below=1.5cm of P1C] (P2D)  {$D$};
    \node[recon, right=of P2D]       (P2T)  {$T'$};
    \node[recon, right=of P2T]       (P2C)  {$C'$};
    \node[metric,right=of P2C]       (P2M)
      {reference-test\\pass rate $R(C', T)$};
    \node[judge, right=of P2M]       (P2J)
      {judge SLM rating\\$J(C, C')$};

    \draw[slm] (P2D) -- node[above, font=\scriptsize]{$L_{\mathsf{test}}$}  (P2T);
    \draw[slm] (P2T) -- node[above, font=\scriptsize]{$L_{\mathsf{code}}$}  (P2C);
    \draw[slm] (P2C) -- (P2M);
    \draw[slm] (P2M) -- (P2J);
    \node[pathlabel, left=0.15cm of P2D]{Path 2};

    %% ------------------- Path 3: C -> T -> D -------------------
    \node[orig,  below=1.5cm of P2D] (P3C)  {$C$};
    \node[recon, right=of P3C]       (P3T)  {$T'$};
    \node[recon, right=of P3T]       (P3D)  {$D'$};
    \node[metric,right=of P3D]       (P3M)
      {BERTScore F1\\$B(D, D')$};
    \node[judge, right=of P3M]       (P3J)
      {judge SLM rating\\$J(D, D')$};

    \draw[slm] (P3C) -- node[above, font=\scriptsize]{$L_{\mathsf{test}}$}  (P3T);
    \draw[slm] (P3T) -- node[above, font=\scriptsize]{$L_{\mathsf{spec}}$}  (P3D);
    \draw[slm] (P3D) -- (P3M);
    \draw[slm] (P3M) -- (P3J);
    \node[pathlabel, left=0.15cm of P3C]{Path 3};

    %% -- Strict-AND validity decision (rounded box below) -------
    \node[draw=black!70, thick, dashed, rounded corners, fill=black!3,
          below=1.0cm of P3M, minimum width=13.5cm, minimum height=0.9cm, font=\small,
          align=center] (validity)
      {\textbf{Strict-AND validity decision (Algorithm~\ref{alg:validity})}:\;
       $\text{valid} \iff (\text{metric} > \tau) \;\wedge\; (\text{judge rating} \geq \rho)$};

    \draw[->, dashed, gray!70] (P1J.south) -- ++(0,-0.15) |- (validity.east);
    \draw[->, dashed, gray!70] (P2J.south) -- ++(0,-0.15) |- (validity.east);
    \draw[->, dashed, gray!70] (P3J.south) -- (validity.north);

    %% -- Legend -----------------------------------------------------
    \node[draw, fill=white, font=\scriptsize, below=0.4cm of validity,
          text width=13cm, align=left, inner sep=4pt]
      {Grey nodes: original artefacts (input). Blue nodes: reconstructed artefacts. Blue box:
       SE-validated automated closure metric. Orange box: external judge SLM rating on
       $\{-1, 0, 1, 2, 3, 4\}$ (DeepSeek-R1:14B, family-disjoint). $L_{\mathsf{spec}}$,
       $L_{\mathsf{test}}$, $L_{\mathsf{code}}$ are the three pipeline SLMs assigned by
       the DOE cell (six candidates; see Table~\ref{tab:doe_summary}). Thresholds are the
       paper defaults $\tau = 0$ and $\rho = 3$.};
  \end{tikzpicture}
  \caption{Round-trip closure framework. Each closure path takes two of the three
    triangle artefacts ($C$: code, $D$: docstring, $T$: reference tests) as
    input and reconstructs the third via two SLM stages. The automated closure
    metric (Path 1: mutation kill rate; Path 2: reference-test pass rate; Path 3:
    BERTScore F1) is paired with an external judge SLM's semantic-equivalence
    rating under the strict-AND validity policy (Algorithm~\ref{alg:validity}).
    Different DOE cells (Section~\ref{sec:experiment_setup}) instantiate
    $L_{\mathsf{spec}}$, $L_{\mathsf{test}}$, $L_{\mathsf{code}}$ with the same
    or different open-weight SLMs.}
  \label{fig:framework}
\end{figure*}


\subsection{Notation and problem setup}
\label{sec:framework_notation}

Table~\ref{tab:notation} summarises the notation used throughout this
paper. In brief: uppercase italic letters $C, D, T$ denote the original
Python code, its docstring, and its reference test suite; primed
variants $C', D', T'$ denote SLM-reconstructed artefacts; $L_{\mathsf{spec}},
L_{\mathsf{test}}, L_{\mathsf{code}}$ denote the pipeline SLMs assigned to
the three stages by the current DOE cell; $L_J$ denotes the external
judge SLM. A closure operator is a two-step composition of two of these
SLMs; the pipeline is \emph{closed} on a given (path, sample) pair when
the reconstruction is semantically equivalent to the original.

\begin{table}[h]
\centering
\small
\caption{Notation used in the round-trip closure formalism.}
\label{tab:notation}
\begin{tabular}{ll}
\toprule
Symbol & Meaning \\
\midrule
$C, D, T$ & original code, docstring, and reference test suite \\
$C', D', T'$ & pipeline-reconstructed artefacts \\
$L_{\textsf{spec}}, L_{\textsf{test}}, L_{\textsf{code}}$ & Small Language Models assigned to the three stages \\
$L_J$ & external judge SLM (DeepSeek-R1, family-disjoint from pipeline) \\
$\mathcal{M} = \{m_1, \ldots, m_6\}$ & the six pipeline SLMs (see Table~\ref{tab:model_lineup}) \\
$\mathcal{S} = \{\textsf{spec}, \textsf{test}, \textsf{code}\}$ & pipeline stages \\
$a : \mathcal{S} \to \mathcal{M} \cup \{\bot\}$ & cell assignment (mapping stages to models or ablation) \\
$K(T', C)$ & mutation kill rate of test suite $T'$ against code $C$ \\
$R(C', T)$ & reference test pass rate of code $C'$ against tests $T$ \\
$B(D_1, D_2)$ & rescaled BERTScore-F1 between two docstrings \\
$J(x, y) \in \{-1, 0, 1, 2, 3, 4\}$ & judge equivalence rating; $-1$ denotes parse failure \\
$\tau$ & metric-validity threshold (path-specific; $\tau = 0$ in our defaults) \\
$\rho$ & judge-validity threshold ($\rho = 3$: ``equivalent or better'') \\
$\mathds{1}[\cdot]$ & indicator function \\
$N$ & per-cell sample size ($150$ for mono, $75$ for hetero and null) \\
\bottomrule
\end{tabular}
\end{table}

\subsection{Closure paths}

We define three closure paths, each traversing two edges of the docstring--test--code triangle. Path 1 (Code $\to$ Docstring $\to$ Tests) tests whether tests derived from an LLM-generated docstring can still detect defects in the original code. Path 2 (Docstring $\to$ Tests $\to$ Code) tests whether code reconstructed from the original docstring and LLM-generated tests can satisfy the reference test suite. Path 3 (Code $\to$ Tests $\to$ Docstring) tests whether the docstring recovered from tests-generated-from-code still describes the original function. Formally,

\begin{align}
  \text{Path 1 (}C \to D \to T\text{):} \quad
  T' &\;=\; L_{\textsf{test}}\bigl(L_{\textsf{spec}}(C)\bigr), \label{eq:path1}\\
  \text{Path 2 (}D \to T \to C\text{):} \quad
  C' &\;=\; L_{\textsf{code}}\bigl(D,\; L_{\textsf{test}}(D)\bigr), \label{eq:path2}\\
  \text{Path 3 (}C \to T \to D\text{):} \quad
  D' &\;=\; L_{\textsf{spec}}\bigl(L_{\textsf{test}}(C)\bigr). \label{eq:path3}
\end{align}

Path 2 explicitly conditions $L_{\textsf{code}}$ on both the original docstring $D$ and the freshly-generated tests $T'$ --- this matches the strongest reconstruction context available to the pipeline. Algorithm~\ref{alg:path_execution} presents the unified execution procedure.

\begin{algorithm}[h]
\caption{Round-trip closure path execution}
\label{alg:path_execution}
\begin{algorithmic}[1]
\Require cell $a = (L_{\textsf{spec}}, L_{\textsf{test}}, L_{\textsf{code}})$, sample $s = (C, D, T)$
\Ensure closure result $r_p$ for each path $p \in \{1, 2, 3\}$
\For{$p \in \{1, 2, 3\}$}
  \If{$p = 1$} \Comment{$C \to D \to T$}
    \State $D' \leftarrow L_{\textsf{spec}}(C)$
    \State $T' \leftarrow L_{\textsf{test}}(D', \textsf{entry\_point}, \textsf{signature})$
    \State $T'_{\textsf{filt}} \leftarrow \textsf{TestFilter}(T', C)$ \Comment{Algorithm~\ref{alg:test_filter}}
    \State $K \leftarrow \textsf{MutationKillRate}(T'_{\textsf{filt}}, C)$
    \State $J \leftarrow \textsf{Judge}(D, D', \text{kind}=\textsc{docstring})$
    \State $r_1 \leftarrow (\text{metric}=K,\; \text{judge}=J)$
  \ElsIf{$p = 2$} \Comment{$D \to T \to C$}
    \State $T' \leftarrow L_{\textsf{test}}(D, \textsf{entry\_point}, \textsf{signature})$
    \State $C' \leftarrow L_{\textsf{code}}(D, T', \textsf{entry\_point}, \textsf{signature})$
    \State $R \leftarrow \textsf{ReferencePassRate}(T, C')$
    \State $J \leftarrow \textsf{Judge}(C, C', \text{kind}=\textsc{code})$
    \State $r_2 \leftarrow (\text{metric}=R,\; \text{judge}=J)$
  \Else \Comment{$p = 3$, $C \to T \to D$}
    \State $T' \leftarrow L_{\textsf{test}}(C)$
    \State $D' \leftarrow L_{\textsf{spec}}(T')$
    \State $B \leftarrow \textsf{BERTScore}(D, D')$
    \State $J \leftarrow \textsf{Judge}(D, D', \text{kind}=\textsc{docstring})$
    \State $r_3 \leftarrow (\text{metric}=B,\; \text{judge}=J)$
  \EndIf
\EndFor
\State \Return $[r_1, r_2, r_3]$
\end{algorithmic}
\end{algorithm}

\subsection{Closure validity decision}

A round-trip closure is considered \emph{valid} when both the automated metric and the judge SLM agree on the artefact equivalence. Formally, for Path 1 we require

\begin{equation}
  \text{valid}_1(C, D, D', T')
  \;=\;
  \mathds{1}\!\bigl[\, K(T', C) > \tau \,\bigr]
  \;\wedge\;
  \mathds{1}\!\bigl[\, J(D, D') \geq \rho \,\bigr],
  \label{eq:valid1}
\end{equation}

and analogously for Paths 2 and 3:

\begin{equation}
  \text{valid}_2 = \mathds{1}[R(C', T) > 0] \wedge \mathds{1}[J(C, C') \geq \rho],
  \label{eq:valid2}
\end{equation}
\begin{equation}
  \text{valid}_3 = \mathds{1}[B(D, D') > 0] \wedge \mathds{1}[J(D, D') \geq \rho].
  \label{eq:valid3}
\end{equation}

The two-signal \emph{strict-AND} policy is deliberate: it prevents metric noise or judge parse failures from spuriously inflating closure rates. We adopt the strict policy because the alternative --- single-signal validity --- was the failure mode responsible for inflated closure rates observed in the 30-function pilot; the strict policy also makes the false-closure rate (Section~5) well-defined.

Every closure row is additionally categorised by Algorithm~\ref{alg:validity} into one of five decision-reason labels that exhaust the two-signal decision space. The disagreement categories (\textsf{false\_closure\_candidate} and \textsf{metric\_false\_negative}) do not count as valid closures, but they become the analytical hooks in Section~5 for characterising \emph{where} the automated metric and the judge diverge.

\begin{algorithm}[h]
\caption{Closure validity decision (implements Eq.~\ref{eq:valid1}--\ref{eq:valid3})}
\label{alg:validity}
\begin{algorithmic}[1]
\Require metric value $m$, judge rating $J \in \{-1, 0, 1, 2, 3, 4\}$, thresholds $\tau$, $\rho$
\Ensure validity $\in \{\text{valid}, \text{invalid}\}$, \textsf{decision\_reason} $\in$ five categories
\If{$m$ is missing \textbf{or} $J = -1$}
  \State \Return $(\text{invalid},\; \textsf{structural\_NA})$ \Comment{Ablation cells or empty artefacts}
\EndIf
\If{$m > \tau$ \textbf{and} $J \geq \rho$}
  \State \Return $(\text{valid},\; \textsf{both\_agree\_valid})$
\ElsIf{$m \leq \tau$ \textbf{and} $J < \rho$}
  \State \Return $(\text{invalid},\; \textsf{both\_agree\_invalid})$
\ElsIf{$m > \tau$ \textbf{and} $J < \rho$}
  \State \Return $(\text{invalid},\; \textsf{false\_closure\_candidate})$ \Comment{Metric over-credits}
\Else \Comment{$m \leq \tau$ and $J \geq \rho$}
  \State \Return $(\text{invalid},\; \textsf{metric\_false\_negative})$ \Comment{Metric under-credits}
\EndIf
\end{algorithmic}
\end{algorithm}

Algorithm~\ref{alg:validity} is implemented as the pure function \texttt{closure\_decision.decide\_validity} in the replication package, with 43 unit tests covering all five decision reasons, the structural-NA edge cases (NaN metric, judge $= -1$, \texttt{None}), and the boundary values of $\tau$ and $\rho$.

\subsection{Test-filter validity gate}

Before computing the mutation kill rate, we discard generated tests that fail on the original code $C$. This avoids penalising mutants for surviving tests that were spuriously wrong to begin with. Formally, given a candidate test set $T' = \{t_1, t_2, \ldots\}$, the filtered suite is

\begin{equation}
  T'_{\textsf{filt}}(C)
  \;=\;
  \bigl\{\, t \in T' \;:\; \textsf{exec}(t, C) = \textsf{pass} \,\bigr\}.
  \label{eq:filter}
\end{equation}

The filter is implemented as Algorithm~\ref{alg:test_filter}, which additionally emits a machine-readable \textsf{filter\_reason} label recording \emph{why} the filter output was empty when it was. This label appears in every closure row in the results TSV and is used in Section~5 to distinguish rows where the LLM produced no tests at all (\textsf{empty\_input}) from rows where every generated test failed on the ground truth (\textsf{all\_dropped}) --- a distinction that matters for the phi-4-in-test-stage analysis in Section~\ref{sec:discussion_rq3}.

\begin{algorithm}[h]
\caption{Test-filter validity gate (implements Eq.~\ref{eq:filter})}
\label{alg:test_filter}
\begin{algorithmic}[1]
\Require candidate test suite $T'$, original code $C$
\Ensure filtered test suite $T'_{\textsf{filt}} \subseteq T'$, \textsf{filter\_reason} label
\If{$T'$ or $C$ is empty}
  \State \Return $(\emptyset,\; \textsf{empty\_input})$
\EndIf
\State $N \leftarrow$ number of \texttt{def test\_*} functions in $T'$
\If{$N = 0$}
  \State \Return $(\emptyset,\; \textsf{no\_test\_functions})$
\EndIf
\State $T'_{\textsf{filt}} \leftarrow \emptyset$
\For{each test $t$ in $\textsf{split\_test\_functions}(T')$}
  \State $\textsf{status} \leftarrow \textsf{pytest\_subprocess}(t, C, \textsf{timeout} = 10\text{s})$
  \If{$\textsf{status} = \textsf{pass}$}
    \State $T'_{\textsf{filt}} \leftarrow T'_{\textsf{filt}} \cup \{t\}$
  \EndIf
\EndFor
\If{$|T'_{\textsf{filt}}| = 0$}
  \State \Return $(\emptyset,\; \textsf{all\_dropped})$ \Comment{Row will be written with NaN metric}
\EndIf
\State \Return $(T'_{\textsf{filt}},\; \textsf{kept}\_K\_\textsf{of}\_N)$
\end{algorithmic}
\end{algorithm}

\subsection{Closure metrics}

\subsubsection{Path 1: mutation kill rate}

Given the filtered test suite $T'_{\textsf{filt}}$ and the set of mutants $\mathcal{M}_C = \{m_1, \ldots, m_n\}$ generated from $C$, the mutation kill rate is

\begin{equation}
  K(T', C)
  \;=\;
  \frac{\bigl|\{\, m \in \mathcal{M}\emph{C \;:\; \exists t \in T'}{\textsf{filt}}(C),\; \textsf{exec}(t, m) = \textsf{fail} \,\}\bigr|}{\bigl|\mathcal{M}_C\bigr| \;-\; \bigl|\mathcal{M}_C^{\textsf{equiv}}\bigr|}.
  \label{eq:killrate}
\end{equation}

A mutant is \emph{killed} if at least one filtered test fails on it. Equivalent mutants $\mathcal{M}_C^{\textsf{equiv}}$ (mutants whose observable behaviour matches the original) are excluded from the denominator following the second companion paper's protocol. Five mutation operators are applied: arithmetic (e.g. $\texttt{+} \to \texttt{-}$), boundary ($\texttt{<} \to \texttt{<=}$), comparison ($\texttt{<} \to \texttt{>}$), negate\_bool, and return\_none. Up to 15 mutants are generated per function.

\subsubsection{Path 2: reference test pass rate}

Let $T = \{t_1, \ldots, t_k\}$ be the ground-truth reference suite. The pass rate is

\begin{equation}
  R(C', T)
  \;=\;
  \frac{1}{|T|}\sum_{i=1}^{|T|} \mathds{1}\!\bigl[\,\textsf{exec}(t_i, C') = \textsf{pass}\,\bigr].
  \label{eq:passrate}
\end{equation}

In this experiment we use a binary indicator ($R \in \{0, 1\}$: the suite either fully passes or does not) for the Section~4 statistical analysis.

\subsubsection{Path 3: BERTScore F1}

For two docstrings $D_1$ and $D_2$, let $\phi(D)$ denote their RoBERTa-large contextual embeddings. Precision and recall under cosine similarity are

\begin{align}
  P(D_1, D_2) &\;=\; \tfrac{1}{|\phi(D_2)|} \textstyle\sum_{x \in \phi(D_2)} \max_{y \in \phi(D_1)} \cos(x, y), \label{eq:bert_p}\\
  R(D_1, D_2) &\;=\; \tfrac{1}{|\phi(D_1)|} \textstyle\sum_{y \in \phi(D_1)} \max_{x \in \phi(D_2)} \cos(x, y), \label{eq:bert_r}
\end{align}

and the BERTScore F1 we report is

\begin{equation}
  B(D_1, D_2)
  \;=\;
  \frac{2 \cdot P(D_1, D_2) \cdot R(D_1, D_2)}{P(D_1, D_2) + R(D_1, D_2)}.
  \label{eq:bertscore}
\end{equation}

We use baseline-rescaled scores so $B \in (-\infty, 1]$ with $B = 0$ denoting ``no better than random pairing''. In our data, $B$ values on real closure pairs concentrate in $[-0.2, 0.7]$.

\subsubsection{Judge equivalence rating}

The judge SLM $L_J$ = DeepSeek-R1:14B is prompted with a 5-level behaviourally-anchored rubric: 4 = identical, 3 = equivalent (same observable behaviour), 2 = approximately equivalent, 1 = clearly different, 0 = unrelated. A rating of $-1$ denotes parse failure --- the judge is treated as absent for that row. DeepSeek-R1 is deliberately chosen because it belongs to a family (DeepSeek) disjoint from every pipeline SLM, avoiding self-evaluation bias.

